\begin{abstract}
Min-sum is the version of belief propagation that is used for solving constraint optimization problems (COPs). Belief propagation is a message passing algorithm in which nodes of the constraint graph (a.k.a.\ factor-graph) take an active role in the algorithm, i.e., they perform calculations and exchange messages with their neighbors. Thus, it is a leading candidate for solving distributed COPs (DCOPs). Unfortunately, the standard vanilla version of Min-sum fails to converge and produces low quality solutions. To mitigate this, previous studies proposed the use of message damping, which encourages the convergence of the algorithm and improves its solution quality. However, damped Min-sum (DMS) requires thousands of iterations in order to converge, and therefore it cannot be used when rapid decision making is required.
Recently, empirical evidence indicated that by splitting the functions of the factor graph on which DMS performs, the algorithm converges immensely faster (almost instantly) to high quality solutions. However, while this success was empirically validated, a theoretical understanding of this phenomenon has not yet been established.

We present theoretical and empirical results that shed light on the success of function-node splitting in improving DMS. We prove that the small cycles generated by such splitting create feedback loops that rapidly separate the solution to which the algorithm converges from alternative candidates, and we establish explicit bounds within which the variable assignments cannot change, i.e., within which the algorithm converges --- including under bounded, dynamically changing influence from the rest of the factor graph. This understanding allows us to propose new heuristics that further improve the algorithm on structured benchmarks, where they significantly outperform a state-of-the-art learned belief-propagation solver.
\end{abstract}
% AAAI-27 \links block: must sit BETWEEN abstract and main body.
% TODO(authors): replace placeholder with the actual anonymised repo URL
% (code + Lean proofs) before submission, and make sure the repo is
% populated. AAAI warns: do not de-anonymise yourself via these links.
\begin{links}
    \link{Code and Lean proofs}{https://anonymized.example/PLACEHOLDER}
\end{links}
